% =============================================================================
% Chapter 5 --- Verification Methodology
% Grounded in Context 3 §1-5 and the actual testbench sources in verif/.
% =============================================================================

\chapter{Verification Methodology}
\label{ch:verification}

This chapter documents the verification environment used in the project.
All claims are grounded in the testbench sources under \sv{verif/}, the
simulation driver \sv{scripts/run\_sim.py}, and the checked-in report
files under \sv{target/}. No UVM environment, constrained-random stimulus
engine, functional coverage model, or formal property proofs are present
in the repository; the verification performed here is directed functional
simulation in ModelSim.

\section{Strategy and Execution Context}
\label{sec:verif-strategy}

The testbenches are written in SystemVerilog, import the three project
packages (\sv{controller\_pkg}, \sv{input\_buffer\_pkg},
\sv{parallel\_interface\_pkg}), and \sv{`include} the shared macro header
\sv{source/common/macros.svh}. Most controller benches instantiate the
three-module stack in the order
\modname{parallel\_interface}\,\(\rightarrow\)\,\modname{ann\_controller}\,\(\rightarrow\)\,\modname{input\_buffer},
matching the intended real connectivity of the system. Unit benches for
\modname{input\_buffer} and \modname{parallel\_interface} instantiate
only the module under test. A block view of this stacking is shown in
\Cref{fig:testbench-stack}.

\begin{figure}[htbp]
    \centering
    \includegraphics[width=0.95\linewidth]{fig_testbench_stack}
    \caption[Testbench stacking]{Testbench stacking for the controller
        benches. The three RTL modules are instantiated in their real
        order; \signame{weight\_read\_data} is driven by a TB-local
        combinational mock that either returns the shadow weight matrix
        cell (baseline) or an injected value (fault injection). The
        \signame{op\_done} handshake is generated by a TB-local
        \sv{always\_ff} that follows the DUT's own sub-FSM state.}
    \label{fig:testbench-stack}
\end{figure}

Each testbench opens its report file with \sv{\$fopen} and writes
pass/fail lines and summaries with \sv{\$fdisplay}/\sv{\$fwrite}. The
reports are closed explicitly and are intended for regression diffing;
they are checked into \sv{target/} and are the primary evidence cited in
\Cref{ch:results}. There is no unified base class or base module; each
testbench is a single self-contained SystemVerilog \sv{module} with local
counters and file handles.

\section{Python-Driven Simulation Flow}
\label{sec:verif-runscript}

The Python helper \sv{scripts/run\_sim.py} wraps \sv{vlog} for compilation
and \sv{vsim} for simulation, and supports the subcommands \sv{list},
\sv{compile}, \sv{sim}, \sv{run\_all}, and \sv{clean}. The expected
ModelSim installation location is documented inside the script as
\sv{/c/intelFPGA\_lite/17.0/modelsim\_ase/win32aloem}. A typical
regression run proceeds as

\begin{lstlisting}[style=plainstyle,language={}]
python scripts/run_sim.py compile -m Controller        -t verif
python scripts/run_sim.py compile -m Input_Buffer      -t verif
python scripts/run_sim.py compile -m Parallel_Interface -t verif
python scripts/run_sim.py run_all -m Controller
python scripts/run_sim.py run_all -m Input_Buffer
python scripts/run_sim.py run_all -m Parallel_Interface
\end{lstlisting}

Individual tests can be selected with \sv{sim -tb <name>}; providing a
\sv{--do-file} switches the run into GUI mode so that the waveform
wrapper testbenches under \sv{verif/.../tb/waves/} can be used with
their companion \sv{.do} scripts. Cleanup operations (\sv{clean -a},
\sv{-w}, \sv{-t}) remove the \sv{work/} library and the
per-module \sv{target/} directories as specified by the script's option
table.

\section{Memristor Behavioral Mock}
\label{sec:verif-mock}

Several controller testbenches --- specifically
\sv{controller\_prog\_verify\_lut\_tb},
\sv{controller\_prog\_verify\_lut\_10w\_tb},
\sv{controller\_host\_erase\_tb}, and
\sv{controller\_host\_read\_reorder\_tb} --- do not instantiate a
separate ANN or ADC model. They maintain a \emph{testbench-local shadow
weight matrix} and drive the DUT input \signame{weight\_read\_data}
from combinational logic. The structure of this mock is described below.

\subsection{Shadow Storage: \texttt{ann\_weight\_matrix}}
\label{sec:verif-mock-matrix}

Each relevant TB declares

\begin{lstlisting}[style=svstyle]
logic [3:0] ann_weight_matrix
    [0:NUM_BLOCKS-1]
    [0:NUM_SUB_BLOCKS-1]
    [0:SUB_BLOCK_ROWS-1]
    [0:SUB_BLOCK_COLS-1];
\end{lstlisting}

Indices follow the same \sv{{block, sub\_block, row, col}} decomposition
as the RTL address map, with the \sv{NUM\_*} constants imported from
\sv{controller\_pkg}. On the negative edge of \signame{rst\_n}, nested
\sv{for} loops clear every element to \sv{4'b0}. The mock is updated
during programming using a shared pattern:

\begin{lstlisting}[style=svstyle]
assign in_prog_core_phase =
    (pulses == PULSE_MODE_PROG) && (buf_reg_ctrl == CTRL_WEIGHT_READ);

// ...

always_ff @(posedge clk) if (in_prog_core_phase) begin
    ann_core_word_decode(ann_core_word, lb, lsb, lr, lc);
    ann_weight_matrix[lb][lsb][lr][lc] <= ann_core_word[27:24];
end
\end{lstlisting}

That is, the nibble programmed onto the core bus (bits
\sv{[27:24]} of \signame{ann\_core\_word}) is what the mock stores at
the decoded cell, mirroring the RTL's own word layout during
programming.

\subsection{Baseline Readback: \texttt{actual\_from\_ann}}
\label{sec:verif-mock-readback}

A combinational block continuously decodes the \emph{current}
\signame{ann\_core\_word} to select the cell whose value the mock
should expose:

\begin{lstlisting}[style=svstyle]
always_comb begin
    ann_core_word_decode(ann_core_word, dec_blk, dec_sb, dec_row, dec_col);
    actual_from_ann = ann_weight_matrix[dec_blk][dec_sb][dec_row][dec_col];
end
\end{lstlisting}

By default, \signame{weight\_read\_data\_mock} is tied to
\signame{actual\_from\_ann} (or to it before optional fault-injection
overrides), and the DUT input is wired as
\sv{.weight\_read\_data(weight\_read\_data\_mock)}. During verify, the
RTL therefore sees either the shadow cell contents or the injected
value described in \Cref{sec:verif-injection}.

Not every testbench uses the shadow matrix.
\sv{controller\_addr\_pulse\_tb} and
\sv{parallel\_interface\_controller\_integration\_tb} tie
\signame{weight\_read\_data} directly to \signame{buf\_data}[3:0], so
verify and read paths see the buffer nibble rather than the shadow
matrix. \sv{controller\_inf\_buffer\_flow\_tb} ties
\signame{weight\_read\_data\_mock} to \signame{buf\_data}[3:0] and
holds \signame{op\_done} low throughout; the INF flow does not exercise
the program--verify mock at all.

\subsection{\texttt{op\_done} Behavioral Mock}
\label{sec:verif-mock-opdone}

The RTL waits on \signame{op\_done} in \sv{PROG\_SELECT},
\sv{PROG\_WAIT\_ACK}, \sv{ERASE\_SELECT}, and \sv{ERASE\_WAIT\_ACK}. The
program--verify--LUT benches model this handshake with an
\sv{always\_ff} that:

\begin{itemize}
    \item Asserts \signame{op\_done}~=~1 when
          \signame{dut.state}~=~\stateval{S\_PROGRAM} and
          \signame{dut.prog\_state} is \sv{PROG\_SELECT} or
          \sv{PROG\_WAIT\_ACK}, or when
          \signame{dut.state}~=~\stateval{S\_ERASE} and
          \signame{dut.erase\_state} is \sv{ERASE\_SELECT} or
          \sv{ERASE\_WAIT\_ACK}.
    \item For other busy intervals in which \signame{pulses} takes
          \sv{PULSE\_MODE\_READ}, \sv{PULSE\_MODE\_ERASE}, or
          \sv{PULSE\_MODE\_INF} (in some benches), a counter
          \signame{op\_done\_cnt} fires a delayed one-cycle pulse so
          that the DUT does not hang.
\end{itemize}

This is testbench-only timing: it is not modeling a specific analog
latency, only enough to let the FSM advance in simulation.

\subsection{End-of-Run Scoreboard against \texttt{weight\_matrix}}
\label{sec:verif-mock-scoreboard}

\sv{controller\_prog\_verify\_lut\_tb} loads expected nibbles from
\sv{target/Controller/programming\_inputs/weight\_matrix.txt} into a
flat array \sv{weight\_matrix[0:639]} and builds per-weight addresses
\sv{weight\_addresses[]} with \sv{matrix\_coords\_to\_address}. After
all \sv{CMD\_PROG} transactions complete, the bench iterates over every
row/column, calls \sv{parse\_ann\_address} to index the shadow array,
and increments an \sv{errs} counter on any
\sv{actual~!==~expected}. The report concludes with a summary line of
the form
\sv{// Summary: \%0d errors in 640 weights}.

\section{Fault Injection for Recovery Coverage}
\label{sec:verif-injection}

A well-behaved mock that returns the exact programmed nibble on every
verify read would always satisfy the comparator in \sv{VERIFY\_CHECK}
on the first attempt, so the recovery branches of the FSM (re-PROG,
verify-initiated erase) would never be exercised. To cover those
branches, the program--verify LUT benches explicitly \emph{inject}
fault values into the verify read window. The injection strategy is
summarized in \Cref{fig:fault-injection-window}.

\begin{figure}[htbp]
    \centering
    \includegraphics[width=0.95\linewidth]{fig_fault_injection_window}
    \caption[Fault injection gating]{Fault injection is active only
        during the verify read window. The bench tracks
        \signame{current\_weight\_idx} across the sweep,
        \signame{in\_verify\_phase} (asserted while the DUT is busy
        and \signame{pulses} is \sv{PULSE\_MODE\_READ} or
        \sv{PULSE\_MODE\_HIZ}), and a small counter
        \signame{verify\_cycle\_cnt}. Together they gate the injection
        to the early verify cycles of the currently-active
        transaction.}
    \label{fig:fault-injection-window}
\end{figure}

The gating signals are:

\begin{itemize}
    \item \signame{current\_weight\_idx} is set at the start of the
          \sv{send\_prog\_and\_wait} task so the combinational override
          knows which weight in the sweep is active.
    \item \signame{in\_verify\_phase}~=~\signame{busy}~\sv{\&\&}~(\signame{pulses}~=~\sv{3'b001}
          \(\lor\) \signame{pulses}~=~\sv{3'b000}), that is,
          the DUT is busy during the verify read train or the
          adjacent low phase where the RTL's \sv{pulses} combinational
          block drives \sv{PULSE\_MODE\_HIZ}.
    \item \signame{verify\_cycle\_cnt} increments once on entry into
          the verify read phase (\signame{busy}~\sv{\&\&}~\signame{pulses}~=~\sv{3'b001}
          \sv{\&\&}~\sv{!was\_in\_verify}) and clears when
          \signame{busy} deasserts.
\end{itemize}

The 640-weight TB activates injection with
\sv{if (current\_weight\_idx~>=~0 \&\& in\_verify\_phase \&\& verify\_cycle\_cnt~<=~1)}.
Combined with the counter behavior, this keeps injection active through
the early verify cycles (READ / WAIT / CHECK) of the current command.

\subsection{Under-Programmed Injection (\texttt{read < expected})}
\label{sec:verif-injection-under}

The bench function \sv{is\_inject\_read\_lt(idx)} is true for the 14
weight indices

\[
    \{5, 15, 25, 50, 100, 200, 385, 420, 455, 490, 535, 580, 605, 625\}.
\]

For those indices, during the gated window, the TB drives
\sv{weight\_read\_data\_mock = weight\_matrix[current\_weight\_idx] - 1}
(with the precondition \sv{weight\_matrix[...]} \(>\) 0). The DUT therefore
sees a readback one LSB step below the expected nibble loaded from the
file, so \sv{VERIFY\_CHECK} takes the under-programmed branch
(\Cref{sec:rtl-controller-verify-check}): while
\signame{prog\_retry\_cnt}~\(<\)~\sv{MAX\_PROG\_RETRIES} (= 3), it sets
\signame{next\_state}~=~\stateval{S\_PROGRAM}; if the retries are
exhausted, it sets \signame{next\_state}~=~\stateval{S\_ERASE}.

\subsection{Over-Programmed Injection (\texttt{read > expected})}
\label{sec:verif-injection-over}

The function \sv{is\_inject\_read\_gt(idx)} is true for the 14 indices

\[
    \{10, 30, 70, 150, 250, 350, 395, 440, 475, 510, 555, 600, 615, 635\}.
\]

For those indices the TB drives
\sv{weight\_read\_data\_mock = weight\_matrix[current\_weight\_idx] + 1}
(guarded so the 4-bit value does not wrap). The DUT sees a readback one
step above the expected nibble, so \sv{VERIFY\_CHECK} asserts the
verify-failure erase branch
(\signame{verify\_failure\_starts\_erase}~=~1,
\signame{next\_state}~=~\stateval{S\_ERASE}). After \sv{ERASE\_COMPLETE}
on the non-host path, if \signame{retry\_cnt}~\(<\)~3, the main state
returns to \stateval{S\_PROGRAM} and programming is re-attempted; the
mock array may be updated again on \signame{in\_prog\_core\_phase} when
the cell is rewritten.

The bench asserts at the end of the run that both
\signame{erase\_phase\_count} and \signame{reprog\_retry\_count} are
nonzero:

\begin{lstlisting}[style=svstyle]
if (erase_phase_count  == 0) $error("Expected ERASE path ...");
if (reprog_retry_count == 0) $error("Expected re-PROG path ...");
\end{lstlisting}

so the regression explicitly fails unless both recovery paths are
exercised.

\subsection{Phase Logging}
\label{sec:verif-injection-phases}

The always block that writes
\sv{target/Controller/prog/prog\_verify\_report.txt} maps
\signame{pulses} and \signame{in\_prog\_core\_phase} to phase strings
\sv{PROG}, \sv{REPROG}, \sv{VERIFY}, \sv{ERASE}, \sv{PROG\_PREP},
\sv{COMPUTE}, and \sv{CHECK\_DONE}. Logging is aligned to
pulse-mode observables, so internal substates such as \sv{PROG\_PREP} or
\sv{COMPUTE} update a \signame{last\_phase} register but may not be
printed on every cycle. This makes the report human-readable while still
distinguishing between a first-attempt PROG and a retry (\sv{REPROG}).

\section{Testbench Catalog}
\label{sec:verif-catalog}

\Cref{tab:tb-catalog} lists every testbench delivered with the project,
its scope, and the output artifact it produces under \sv{target/}.

\begin{table}[htbp]
    \centering
    \caption{Testbench catalog. Paths are relative to the repository
        root.}
    \label{tab:tb-catalog}
    \small
    \begin{tabularx}{\linewidth}{l l X}
        \toprule
        \textbf{Testbench} & \textbf{DUT} & \textbf{Report / Purpose} \\
        \midrule
        \sv{controller\_addr\_pulse\_tb} & Controller &
            30 directed tests of address packing and pulse mode;
            \sv{target/Controller/controller\_addr\_pulse\_verify.txt}. \\
        \sv{parallel\_interface\_controller\_integration\_tb} &
            PI + Ctrl + Buf &
            10 mixed READ/PROG/ERASE/INF transactions;
            \sv{target/Controller/tb\_pi\_controller\_integration.txt}. \\
        \sv{controller\_prog\_verify\_lut\_tb} & Controller &
            640-weight PROG+VERIFY with injected under/over failures and
            LUT-based PROG;
            \sv{target/Controller/prog/prog\_verify\_report.txt}. \\
        \sv{controller\_prog\_verify\_lut\_10w\_tb} & Controller &
            Compact 10-weight program--verify for debug;
            \sv{target/Controller/prog/prog\_verify\_10w\_report.txt}. \\
        \sv{controller\_host\_erase\_tb} & Controller &
            Host-initiated \sv{CMD\_ERASE} with retention check;
            \sv{target/Controller/erase/controller\_host\_erase\_report.txt}. \\
        \sv{controller\_host\_read\_reorder\_tb} & Controller &
            8 weights read in permuted order;
            \sv{target/Controller/read/controller\_host\_read\_reorder\_report.txt}. \\
        \sv{controller\_inf\_buffer\_flow\_tb} & Controller &
            8-pixel INF collect \(\rightarrow\) compute
            \(\rightarrow\) bit-serial check;
            \sv{target/Controller/inf/controller\_inf\_buffer\_flow\_tb\_log.txt}. \\
        \sv{input\_buffer\_bit\_serial\_tb} & Input Buffer &
            Bit-serial lane correctness; \sv{target/Input\_Buffer/}. \\
        \sv{input\_buffer\_reset\_behavior\_tb} & Input Buffer &
            Reset-to-zero check; \sv{target/Input\_Buffer/}. \\
        \sv{input\_buffer\_full\_overwrite\_tb} & Input Buffer &
            Full 64-byte overwrite sweep; \sv{target/Input\_Buffer/}. \\
        \sv{parallel\_interface\_extract\_tb} & Parallel Interface &
            Valid/invalid tail decode checks;
            \sv{target/Parallel\_Interface/parallel\_interface\_extract\_tb\_log.txt}. \\
        \bottomrule
    \end{tabularx}
\end{table}

\section{Tasks and Scoreboard Conventions}
\label{sec:verif-tasks}

The testbenches share a handful of recurring task patterns:

\begin{itemize}
    \item \sv{load\_weights} / \sv{load\_weights\_10}: file I/O into
          \sv{weight\_matrix} and address construction via
          \sv{matrix\_coords\_to\_address}.
    \item \sv{send\_prog\_and\_wait}: waits while \signame{busy} is high,
          issues a one-cycle \signame{host\_cmd}~=~\sv{CMD\_PROG} with
          a \sv{build\_host\_ann\_word}-formed payload, and waits for
          \signame{busy} to rise and then fall, with a cycle-based
          timeout.
    \item \sv{run\_cmd\_test} (in \sv{controller\_addr\_pulse\_tb}):
          drives a command, then compares \signame{ann\_core\_word}
          against \signame{expected\_ann\_core\_word}. For
          \sv{CMD\_READ}, it scans until the first match; for others it
          samples after a few cycles.
    \item \sv{check\_cmd} (integration TB): single-beat transaction with
          idle-gated completion on \signame{busy}; see
          \Cref{sec:verif-integration}.
    \item \sv{check\_inf\_row8} (integration TB): a nine-cycle
          \sv{CMD\_INF} burst with a parallel monitor; see
          \Cref{sec:verif-integration}.
    \item \sv{send\_inf\_beat} / \sv{dump\_row0}
          (\sv{controller\_inf\_buffer\_flow\_tb}): per-beat write and
          peek at buffer row 0.
\end{itemize}

Checks are directed: per-test expected values (computed with the same
package functions used by the RTL) are compared to observed RTL outputs
and --- via hierarchical peeks such as
\sv{u\_input\_buffer.buffer\_reg[...]} --- to internal buffer contents.
Failures increment integer counters and emit \sv{FAIL} lines; some
benches call \sv{\$fatal} or \sv{\$error} on timeout or nonzero failure
count.

\section{Integration Testbench}
\label{sec:verif-integration}

\sv{parallel\_interface\_controller\_integration\_tb.sv} is the
integration bench. Its salient structural choices are:

\begin{itemize}
    \item The DUT stack is
          \modname{parallel\_interface}\,\(\rightarrow\)\,\modname{ann\_controller}\,\(\rightarrow\)\,\modname{input\_buffer},
          with \sv{valid}, \sv{pi\_data}, \sv{address}, and \sv{cmd}
          passed as on the real datapath.
    \item Verify readback is tied as
          \sv{.weight\_read\_data(buf\_data[3:0])}, so integration
          checks do not use the shadow matrix.
\end{itemize}

The \sv{check\_cmd} task implements a single-beat host transaction:

\begin{enumerate}
    \item \sv{wait\_until\_idle(maxcyc)} spins while \signame{busy} is
          high, up to a cycle budget.
    \item The host bus is quieted: \signame{host\_cmd}~=~\sv{CMD\_HIZ},
          \signame{host\_data}~=~0, held for a few clock edges.
    \item Stimulus is issued for one cycle with
          \signame{host\_data}~=~\sv{build\_host\_ann\_word(d,a)} and
          \signame{host\_cmd}~=~c, then deasserted. Holding the command
          high for multiple cycles would incorrectly re-arm
          \sv{CMD\_INF}, so the task deliberately drops it after one
          beat for non-INF commands.
    \item \sv{wait(busy)} blocks until the controller accepts the
          transaction.
    \item While \signame{busy}, on each positive clock edge: cycle
          count is incremented; \sv{pulse\_ok} is set once
          \signame{pulses}~=~\signame{exp\_pulse}; \sv{word\_ok} is set
          once \signame{ann\_core\_word}~=~\sv{exp\_word(a,d)}.
    \item Both \sv{pulse\_ok} and \sv{word\_ok} must be true for
          \sv{pass\_c++}; otherwise the bench logs a \sv{FAIL} line.
          A still-busy DUT after the cycle budget is counted as an
          additional failure.
    \item \sv{dump\_buffer\_snapshot} reads
          \sv{u\_input\_buffer.buffer\_reg[0:63]} and prints an 8-by-8
          hex snapshot.
\end{enumerate}

For inference, the bench uses \sv{check\_inf\_row8}, which launches a
\sv{fork}: one branch drives nine consecutive \sv{CMD\_INF} beats with
the same payload; the parallel branch waits for \signame{busy} and then
runs the same monitor loop as \sv{check\_cmd}, this time expecting
\sv{PULSE\_MODE\_INF} at least once and tolerating the fact that the
\signame{ann\_core\_word} may toggle between COLLECT and COMPUTE
phases. The two branches \sv{join} before the buffer dump and idle
check. The overall run drives ten labelled scenarios (\sv{01\_READ},
\sv{02\_PROG}, \ldots, \sv{10\_ERASE}) and prints a summary of the form
\sv{// Summary: PASS=\%0d FAIL=\%0d}, with \sv{\$fatal} called on a
non-zero \sv{fail\_c}.

\section{Unit Testbenches and Wave Wrappers}
\label{sec:verif-unit}

The input buffer is exercised by three unit benches:
\sv{input\_buffer\_bit\_serial\_tb}, \sv{input\_buffer\_reset\_behavior\_tb},
and \sv{input\_buffer\_full\_overwrite\_tb}. Each performs directed
writes and compares against \sv{buffer\_reg}, reporting with
\sv{SUMMARY: N passed, M failed} in the checked-in log.
The parallel interface is covered by
\sv{parallel\_interface\_extract\_tb}, which exercises decode
assumptions using \sv{build\_host\_ann\_word} and related patterns;
the log under \sv{target/Parallel\_Interface/} carries six checks and a
summary line.

\emph{Wave wrappers} (testbenches whose names end in \sv{\_waves\_tb})
instantiate the corresponding regular TB without adding any additional
checks. They exist so that ModelSim GUI runs can load the companion
\sv{.do} scripts under \sv{verif/.../do/waves/} to preload signal groups
(FSM states, pulse counters, LUT fields, \signame{D0}\ldots\signame{D7})
for debug. See the user-guide chapter of the project documentation for
the full list of wave commands; those runs do not contribute new
regression evidence.

\section{Pass/Fail Accounting Summary}
\label{sec:verif-passfail}

Each testbench carries its own counters (most commonly \signame{errs},
\signame{pass\_c}/\signame{fail\_c}, or per-item PASS/ERROR lines) and
emits a final summary. These counters are what \Cref{ch:results}
quotes: the verdict of each bench is a direct read of its report file,
not a re-computation from the checked-in transcripts.
